\section{Lecture 10 (October 14th)}
\begin{thm}
(Arzerla Ascoli) Let $\{f_{n}\}$ be a pointwise bounded and equicontinuous sequence of functions on a separable metric space $X$. Then, $\{f_{n}\}$ has a subsequence which converges uniformly on all compact subsets of $X$. This subsequence converges pointwise since every $\{x_{n}\}\subset X$ is compact. 
\end{thm}
\vspace{2ex}
\begin{defi}
(Pointwise bounded) A sequence $\{f_{n}\}$ is pointwise bounded on $X$ provided that for every $x\in X$ there is $M_{x}>0$ such that $|f_{n}(x)|\leq M_{x}$ for all $n\in {\bm N}$. 
\end{defi}
\vspace{2ex}
\begin{defi}
(Equicontinuity) A sequence $\{f_{n}\}$ is equicontinuous on $X$ provided that for every $\varepsilon >0$, there is $\delta >0$ such that if $d(x,y)<\delta $ then $|f_{n}(x)-f_{n}(y)|<\varepsilon $ for all $n\in {\bm N}$.
\end{defi}
\vspace{2ex}
\begin{rmk}
If $\{f_{n}\}$ is equicontinuous on $X$, then each $f_{n}$ is uniformly continuous on $X$. The converse is not true, for example set $f_{n}:{\bm R}\rightarrow {\bm R}$ as 
\[f_{n}(x)=nx\]
Notice how $|f_{n}(x)-f_{n}(y)|=n|x-y|$ and each $f_{n}$ is uniformly continuous, but $\{f_{n}\}$ is not equicontinuous on ${\bm R}$.
\end{rmk}
\vspace{2ex}
\begin{proof}
Let $E=\{x_1,x_2,\ldots \}$ be a countable dense subset of $X$. The proof is done in two steps:

\bigskip

(Step 1) We need to use pointwise boundedness to find a subsequence $\{g_{n}\}$ of $\{f_{n}\}$ which converges pointwise on $E$. Notice that $\{f_{n}(x_1)\}$ is a bounded sequence in ${\bm C}$, such that it has a convergent subsequence $\{f_{n,1}(x_1)\}$. Then, $\{f_{n,1}(x_2)\}$ is bounded and has a bounded subsequence in ${\bm C}$, such that it has a convergent subsequence $\{f_{n,2}(x_2)\}$. Notice that meanwhile, $\{f_{n,2}(x_1)\}$ also converges. Continuing and taking the diagonal, we have $\{f_{n,n}(x)\}$ which converges for every $x_{n}\in E$. Set $g_{n}=f_{n,n}$, and we see that for any $x_{m}$, $g_{m},g_{m+1},\ldots $ is a subsequence of $\{f_{n,m}\}^{\infty }_{n=1}$ which converges at $x_{m}$. 

\bigskip

(Step 2) We now use equicontinuity to prove that this sequence converges uniformly on any compact subset $K$ of $X$. We now use equicontinuity to show that there exists a $N\in {\bm N}$ satisfying if $n,m>N$ then $|f_{n}(x)-f_{m}(x)|<\varepsilon $ for every $x\in K$ ($\{f_{n}\}$ is uniformly Cauchy on $K$). Observe that due to the equicontinuity of $\{g_{n}\}$, for every $\varepsilon >0$, there exists $\delta >0$ such that if $p,q\in X$ satisfies $d(p,q)<\delta $, then $|f_{n}(p)-f_{n}(q)|<\varepsilon /3$. Also, 
\[\mathcal{F}=\Big\{B\Big(x,\dfrac{\delta }{2}\Big)\;\Big|\;x\in K\Big\}\]
is an open cover of $K$ so that there is a finite subcover $\{B_1,B_2,\ldots ,B_{M}\}$ (here, note that if $x,y\in B_{k}$, then $d(x,y)<\delta $). Since $E$ is dense in $X$, then $p_{i}\in E\cap B_{i}$ for $1\leq i\leq M$. From step 1, $\lim _{n\rightarrow \infty } g_{n}(p_{i})$ converges for $1\leq i\leq M$ and thus $\{g_{n}(p_{i})\}$ is a Cauchy sequence for $1\leq i\leq M$. Thus there is $N\in {\bm N}$ such that if $n,m>N$ then $|g_{n}(p_{i})-g_{m}(p_{i})|<\varepsilon /3$ for $1\leq i\leq M$ with $N=\max\{N_1,\ldots ,N_{M}\}$. Now we show the same for an arbitrary $x\in K$.
\[\begin{align*}
|g_{n}(x)-g_{m}(x)|=&|g_{n}(x)-g_{n}(p_{k})+g_{n}(p_{k})-g_{m}(p_{k})+g_{m}(p_{k})-g_{m}(x)|\\
\leq&|g_{n}(x)-g_{n}(p_{k})|+|g_{n}(p_{k})-g_{m}(p_{k})|+|g_{m}(p_{k})-g_{m}(x)|\\
\leq&\varepsilon 
\end{align*}\]
as $d(x,p_{k}),d(x,p_{m})<\delta $.
\end{proof}
\vspace{2ex}
\begin{defi}
(Uniformly bounded on a compact $K$) For every compact $K\subset D$ , there is $M_{K}>0$ such that $|f_{n}(z)|\leq M_{K}$ for all $z\in K$, $m\in {\bm N}$.
\end{defi}
\vspace{2ex}
\begin{thm}
(Montel's theorem) If $\{f_{n}\}$ is a squence of analytic functions on a domain $D$ which is uniformly bounded on each compact subset of $D$, then $\{f_{n}\}$ has a subsequence which converges uniformly on all compact subsets of $D$.
\end{thm}
\vspace{2ex}
\begin{thm}
If $\{f_{n}\}$ is a sequence of analytic functions on a domain $D$ which is uniformly bounded on a compact subset of $D$, then $\{f_{n}\}$ is equicontinuous on every compact subset of $D$. 
\end{thm}
\vspace{2ex}
\begin{proof}
The proof is done through the Cauchy integral formula along with the exhaustion of $D$ by compact subsets. 
\end{proof}
\vspace{2ex}

